% SEGMENT OLP-0052-S01
% Part: sets-functions-relations
% Chapter: infinite
% Section: dedekind-induction

\documentclass[../../../include/open-logic-section]{subfiles}

\begin{document}

\olfileid{sfr}{infinite}{induction}
\olsection[கணிதத் தொகுத்தறிதல்]{டெடெகிண்ட் இயற்கணிதங்களும் கணிதத் தொகுத்தறிதலும்}

மிக முதன்மையாக, இப்போது ஒரு டெடெகிண்ட் இயற்கணிதம்---உண்மையில்,
\emph{எந்த} டெடெகிண்ட் இயற்கணிதமும்---இயல் எண்களுக்கான பதிலியாகச்
செயல்படும். பின்வரும் நேரடி விளைவே இதற்குக் காரணம்:

% SEGMENT OLP-0052-S02
\begin{thm}[கணிதத் தொகுத்தறிதல்]\ollabel{thm:dedinfiniteinduction}
$N, s, o$ ஆகியவை ஒரு டெடெகிண்ட் இயற்கணிதத்தை அமைக்கட்டும்.
அப்போது எந்தக் கணம் $X$ க்கும்:
	\begin{center}
		$o \in X$ எனவும் $(\forall {n} \in N \cap X){s}({n}) \in X$
        எனவும் இருந்தால், $N \subseteq X$.
	\end{center}
\end{thm}

% SEGMENT OLP-0052-S03
\begin{proof}
டெடெகிண்ட் இயற்கணிதத்தின் வரையறைப்படி,
$N = \closureofunder{s}{o}$.
இப்போது ${o} \in X$ எனவும் $(\forall {n} \in N)(n \in X \lif
{s}({n}) \in X)$ எனவும் இருந்தால்,
$N = \closureofunder{s}{o} \subseteq X$.
\end{proof}

தொகுத்தறிதல் இயல் எண்களின் இயல்புக்கூறாக இருப்பதால், கருத்து இதுதான்.
எந்த டெடெகிண்ட்-முடிவுறாத கணம் கொடுக்கப்பட்டாலும், ஒரு டெடெகிண்ட்
இயற்கணிதத்தை அமைத்து, அந்த இயற்கணிதத்தை இயல் எண்களுக்கான பதிலியாகப்
பயன்படுத்தலாம்.

% SEGMENT OLP-0052-S04
ஒப்புக்கொள்ள வேண்டுமெனில், \olref{thm:dedinfiniteinduction} ஆனது
தொகுத்தறிதலைக் \emph{கணக்கோட்பாட்டு} சொற்களில் வாய்பாடாக்குகிறது.
ஆனால் அதே நெறியை இன்னும் பழக்கமான சொற்களில் எளிதாகக் கூறலாம்:

\begin{cor}\ollabel{natinductionschema}
$N, s, o$ ஆகியவை ஒரு டெடெகிண்ட் இயற்கணிதத்தை அமைக்கட்டும்.
அப்போது அளபுருக்களைக் கொண்டிருக்கக்கூடிய எந்த வாய்பாடு
$\phi(x)$ க்கும்:
\begin{center}
	$\phi(o)$ எனவும் $(\forall {n} \in N)(\phi(n)\lif
	\phi({s}({n})))$ எனவும் இருந்தால், $(\forall n \in N)\phi(n)$.
\end{center}
\end{cor}

% SEGMENT OLP-0052-S05
\begin{proof}
$X = \Setabs{n \in N}{\phi(n)}$ என்க; இப்போது
\olref{thm:dedinfiniteinduction} ஐப் பயன்படுத்துக.
\end{proof}

% SEGMENT OLP-0052-S06
இந்த முடிவில், ஒரு வாய்பாடு ``அளபுருக்களைக் கொண்டிருப்பது'' பற்றிப்
பேசினோம். இதன் தோராயமான பொருள்: எந்தப் பொருள்கள்
$c_1, \ldots, c_k$ க்கும், நாம் $\phi(x, c_1, \ldots, c_k)$ உடன்
வேலை செய்யலாம் என்பதாகும். இன்னும் துல்லியமாக, ``அளபுருக்கள்'' என்று
குறிப்பிடாமலேயே முடிவைப் பின்வருமாறு கூறலாம். கட்டற்ற மாறிகள் அனைத்தும்
காட்டப்பட்டுள்ள எந்த வாய்பாடு $\phi(x, v_1, \ldots, v_k)$ க்கும்:
	\begin{align*}
			\forall v_1 \ldots \forall v_k((&\phi(o, v_1,\ldots, v_k) \land {}\\
			&	(\forall x \in N)(\phi(x,v_1, \ldots, v_k) \lif \phi(s(x), v_1,\ldots, v_k))) \lif {}\\
			&\hspace{3em} (\forall x \in N)\phi(x, v_1,\ldots, v_k))
	\end{align*}
``அளபுருக்களைக் கொண்டிருப்பது'' என்று பேசுவது வாசிப்பை மிகவும்
எளிதாக்கும் என்பது தெளிவு. (\olref[sth][][]{part} இல், இந்த
உத்தியை அடிக்கடி பயன்படுத்துவோம்.)

% SEGMENT OLP-0052-S07
டெடெகிண்ட் இயற்கணிதங்களுக்குத் திரும்புவோம்: எந்த டெடெகிண்ட்
இயற்கணிதம் கொடுக்கப்பட்டாலும், கூட்டல், பெருக்கல், அடுக்கேற்றம் என்ற
வழக்கமான எண்கணிதச் சார்புகளையும் வரையறுக்கலாம். ஆனால் இது
எளியதன்று; \emph{மறுநிகழ்வு வரையறை} என்ற உத்தி இதில் பயன்படுகிறது.
இந்த உத்தியை மிகவும் பின்னர், இன்னும் பொதுவான சூழலில் அறிமுகப்படுத்தி
நியாயப்படுத்துவோம். (ஆர்வமுள்ளோர் \olref[sth][ord-arithmetic][]{chap}
க்குப் பிறகு இதனை மீண்டும் பார்க்க விரும்பலாம்; அல்லது
\citealt[pp.~95--8]{Potter2004} போன்ற மாற்று விளக்கத்தைப் படிக்கலாம்.)
ஆனால் $N, s, o$ ஆகியவை ஒரு டெடெகிண்ட் இயற்கணிதத்தை அமைக்கும்போது,
இறுதியில் பின்வருவனவற்றை விதித்து:
\begin{align*}
	{a} + {o} &= {a} & & & {a} \times {o} &= {o} & & & {a}^{o} &= s(o)\\
{a} + {s}({b}) &= {s}({a}+{b}) &&& {a} \times {s}({b}) &= ({a}\times {b}) + {a}  & & & {a}^{{s}({b})} &= {a}^{b} \times {a}
\end{align*}
அவை எதிர்பார்த்தபடி செயல்படுகின்றன என்று காட்ட முடியும்.

\end{document}
